\documentclass[11pt, a4paper]{article}

% ─────────────────────────────────────────────────
% Paquetes
% ─────────────────────────────────────────────────
\usepackage[utf8]{inputenc}
\usepackage[T1]{fontenc}
\usepackage[spanish]{babel}
\usepackage{lmodern}
\usepackage{geometry}
\usepackage{graphicx}
\usepackage{xcolor}
\usepackage{tikz}
\usetikzlibrary{arrows.meta, positioning, shapes.geometric, fit, calc}
\usepackage{enumitem}
\usepackage{booktabs}
\usepackage{tabularx}
\usepackage{longtable}
\usepackage{hyperref}
\usepackage{fancyhdr}
\usepackage{titlesec}
\usepackage{float}
\usepackage{tcolorbox}
\tcbuselibrary{skins, breakable}

\geometry{
  left=2.5cm, right=2.5cm,
  top=2.8cm, bottom=2.5cm
}

% ─────────────────────────────────────────────────
% Paleta de colores
% ─────────────────────────────────────────────────
\definecolor{primary}{HTML}{1E3A5F}
\definecolor{accent}{HTML}{2196F3}
\definecolor{success}{HTML}{4CAF50}
\definecolor{warning}{HTML}{FFC107}
\definecolor{danger}{HTML}{F44336}
\definecolor{muted}{HTML}{90A4AE}
\definecolor{bglight}{HTML}{F5F7FA}
\definecolor{bgdark}{HTML}{E8EDF2}

% ─────────────────────────────────────────────────
% Estilos de título
% ─────────────────────────────────────────────────
\titleformat{\section}
  {\Large\bfseries\color{primary}}
  {\thesection}{1em}{}
  [\vspace{-0.4em}\textcolor{accent}{\rule{\linewidth}{1.2pt}}]

\titleformat{\subsection}
  {\large\bfseries\color{primary!85}}
  {\thesubsection}{0.8em}{}

\titleformat{\subsubsection}
  {\normalsize\bfseries\color{primary!70}}
  {\thesubsubsection}{0.6em}{}

% ─────────────────────────────────────────────────
% Cajas de información
% ─────────────────────────────────────────────────
\newtcolorbox{reglanegocio}[1][]{
  enhanced, breakable,
  colback=accent!5, colframe=accent!70,
  fonttitle=\bfseries\color{white},
  coltitle=white, colbacktitle=accent!80,
  title={#1},
  left=6pt, right=6pt, top=4pt, bottom=4pt,
  boxrule=0.6pt, arc=3pt
}

\newtcolorbox{restriccion}[1][]{
  enhanced, breakable,
  colback=danger!5, colframe=danger!60,
  fonttitle=\bfseries\color{white},
  coltitle=white, colbacktitle=danger!70,
  title={#1},
  left=6pt, right=6pt, top=4pt, bottom=4pt,
  boxrule=0.6pt, arc=3pt
}

\newtcolorbox{nota}[1][]{
  enhanced, breakable,
  colback=bglight, colframe=muted!60,
  fonttitle=\bfseries,
  title={#1},
  left=6pt, right=6pt, top=4pt, bottom=4pt,
  boxrule=0.5pt, arc=3pt
}

% ─────────────────────────────────────────────────
% Header / Footer
% ─────────────────────────────────────────────────
\pagestyle{fancy}
\fancyhf{}
\fancyhead[L]{\small\textcolor{primary}{\textbf{SGI} — Lógica de Negocio}}
\fancyhead[R]{\small\textcolor{muted}{v1.0 — Julio 2026}}
\fancyfoot[C]{\small\textcolor{muted}{\thepage}}
\renewcommand{\headrulewidth}{0.4pt}
\renewcommand{\footrulewidth}{0pt}

% ─────────────────────────────────────────────────
% Documento
% ─────────────────────────────────────────────────
\begin{document}

% ── Portada ──
\begin{titlepage}
  \centering
  \vspace*{3cm}

  {\Huge\bfseries\textcolor{primary}{Sistema de Gestión de\\[0.3em]Incidencias Georreferenciadas}}

  \vspace{0.8cm}
  {\Large\textcolor{accent}{Documento de Lógica de Negocio}}

  \vspace{0.5cm}
  {\large\textcolor{muted}{v1.0 — Julio 2026}}

  \vspace{2cm}

  \textcolor{primary}{\rule{0.6\textwidth}{1.5pt}}

  \vspace{1.5cm}

  {\large
    \textbf{Arquitectura:} Hexagonal / DDD (Domain-Driven Design)\\[0.4em]
    \textbf{Backend:} Laravel 12 · PHP 8.3 · PostgreSQL/PostGIS\\[0.4em]
    \textbf{Frontend:} HTML + CSS + JavaScript · Bootstrap · AdminLTE
  }

  \vfill
  {\small\textcolor{muted}{Documento generado automáticamente desde el código fuente del proyecto.}}
\end{titlepage}

\tableofcontents
\newpage

% ═══════════════════════════════════════════════════
\section{Visión General del Sistema}
% ═══════════════════════════════════════════════════

SGI es una aplicación web diseñada para el registro, seguimiento, asignación y resolución de incidencias urbanas georreferenciadas. Opera bajo un modelo multi-rol donde ciudadanos reportan problemas de infraestructura pública y equipos operativos los gestionan dentro de zonas territoriales definidas.

\subsection{Módulos del sistema}

El backend está organizado en \textbf{Bounded Contexts} siguiendo DDD:

\begin{table}[H]
\centering
\begin{tabularx}{\textwidth}{l X}
  \toprule
  \textbf{Módulo} & \textbf{Responsabilidad} \\
  \midrule
  \texttt{Auth}              & Autenticación, autorización, roles, permisos, 2FA, identidades OAuth \\
  \texttt{Incidents}         & Ciclo de vida completo de las incidencias \\
  \texttt{Operations}        & Perfiles operativos, asignaciones supervisor--operador, carga de trabajo \\
  \texttt{TerritorialUnits}  & Jerarquía territorial nacional (país → zona → provincia → cantón → parroquia) \\
  \texttt{Catalogs}          & Catálogos maestros: categorías, prioridades, estados \\
  \texttt{Audit}             & Trazabilidad automática de cambios en entidades auditables \\
  \texttt{Users}             & CRUD administrativo de usuarios \\
  \bottomrule
\end{tabularx}
\caption{Bounded Contexts del sistema}
\end{table}

\subsection{Estructura interna de cada módulo}

Cada módulo sigue la arquitectura hexagonal con tres capas:

\begin{enumerate}[leftmargin=2cm]
  \item \textbf{Domain} — Entidades, Value Objects, interfaces de repositorio, reglas de negocio puras.
  \item \textbf{Application} — Casos de uso (Use Cases), DTOs de entrada/salida, puertos de aplicación.
  \item \textbf{Infrastructure} — Implementaciones concretas: Eloquent, controladores HTTP, mappers, notificadores.
\end{enumerate}


% ═══════════════════════════════════════════════════
\section{Usuarios, Roles y Permisos}
% ═══════════════════════════════════════════════════

\subsection{Modelo de usuario}

Un \textbf{usuario} (\texttt{auth.users}) es la entidad central de identidad. Soporta:

\begin{itemize}
  \item Soft deletes (\texttt{deleted\_at}) para eliminación lógica.
  \item Desactivación temporal (\texttt{is\_active = false}).
  \item Verificación de email obligatoria (\texttt{MustVerifyEmail}).
  \item Autenticación por tokens Sanctum (API stateless).
  \item Autenticación 2FA con TOTP (\texttt{two\_factor\_secret}, \texttt{two\_factor\_confirmed\_at}).
  \item Vinculación con proveedores OAuth (Google) via \texttt{auth.user\_identities}.
\end{itemize}

\subsection{Sistema de roles}

El sistema implementa un esquema \textbf{RBAC} (Role-Based Access Control) con un único rol por usuario.

\begin{table}[H]
\centering
\begin{tabularx}{\textwidth}{l l X}
  \toprule
  \textbf{Código} & \textbf{Nombre} & \textbf{Descripción} \\
  \midrule
  \texttt{ADMIN}      & Administrador & Control total del sistema. Todos los permisos. \\
  \texttt{SUPERVISOR}  & Supervisor    & Gestión de zonas, equipo de operadores, aprobación de cambios de estado. \\
  \texttt{OPERADOR}    & Operador      & Resolución de incidencias asignadas en su zona. \\
  \texttt{CIUDADANO}   & Ciudadano     & Reporte de incidencias y seguimiento de las propias. \\
  \bottomrule
\end{tabularx}
\caption{Roles del sistema}
\end{table}

\begin{restriccion}[Restricción: Rol único]
Cada usuario tiene exactamente un rol asignado. La migración \texttt{enforce\_single\_role\_per\_user} garantiza esta restricción a nivel de base de datos con un índice único en \texttt{auth.role\_user(user\_id)}.
\end{restriccion}

\subsection{Permisos}

Los permisos son granulares y se asignan a cada rol. La relación es \texttt{roles $\leftrightarrow$ permissions} (muchos a muchos via \texttt{auth.permission\_role}).

\begin{longtable}{l l p{7.5cm}}
  \toprule
  \textbf{Módulo} & \textbf{Código} & \textbf{Descripción} \\
  \midrule
  \endfirsthead
  \toprule
  \textbf{Módulo} & \textbf{Código} & \textbf{Descripción} \\
  \midrule
  \endhead
  Incidencias & \texttt{incidents.view}    & Consultar listado y detalle \\
              & \texttt{incidents.create}  & Registrar nuevas incidencias \\
              & \texttt{incidents.edit}    & Actualizar datos operativos \\
              & \texttt{incidents.delete}  & Eliminar incidencias \\
              & \texttt{incidents.assign}  & Asignar operadores \\
              & \texttt{incidents.close}   & Marcar como resueltas \\
              & \texttt{incidents.reopen}  & Reabrir incidencias cerradas o rechazadas \\
  \midrule
  Comentarios & \texttt{comments.create}   & Comentarios públicos \\
              & \texttt{comments.internal} & Comentarios internos (solo equipo operativo) \\
  \midrule
  Usuarios    & \texttt{users.view}         & Consultar usuarios \\
              & \texttt{users.create}       & Crear usuarios manualmente \\
              & \texttt{users.edit}         & Editar información \\
              & \texttt{users.delete}       & Desactivar / eliminar \\
              & \texttt{users.manage\_roles} & Gestionar roles y permisos \\
  \midrule
  Operaciones & \texttt{operations.view}       & Ver cobertura operativa \\
              & \texttt{operations.view\_team}  & Ver equipo de trabajo propio \\
              & \texttt{operations.manage}     & Gestionar cobertura operativa \\
  \midrule
  Catálogos   & \texttt{catalogs.manage}       & Administrar catálogos maestros \\
  \midrule
  Auditoría   & \texttt{audit.view}            & Consultar trazabilidad \\
  \bottomrule
  \caption{Permisos principales del sistema}
\end{longtable}

\subsection{Matriz de permisos por rol}

\begin{table}[H]
\centering
\small
\begin{tabular}{l c c c c}
  \toprule
  \textbf{Permiso} & \textbf{Admin} & \textbf{Supervisor} & \textbf{Operador} & \textbf{Ciudadano} \\
  \midrule
  incidents.view      & \checkmark & \checkmark & \checkmark & \checkmark \\
  incidents.create    & \checkmark & ---        & ---        & \checkmark \\
  incidents.edit      & \checkmark & \checkmark & \checkmark & --- \\
  incidents.assign    & \checkmark & \checkmark & ---        & --- \\
  incidents.close     & \checkmark & \checkmark & ---        & --- \\
  incidents.reopen    & \checkmark & \checkmark & ---        & --- \\
  comments.internal   & \checkmark & \checkmark & \checkmark & --- \\
  users.manage\_roles & \checkmark & ---        & ---        & --- \\
  operations.manage   & \checkmark & ---        & ---        & --- \\
  audit.view          & \checkmark & ---        & ---        & --- \\
  \bottomrule
\end{tabular}
\caption{Permisos clave por rol (resumen)}
\end{table}

\subsection{Autenticación de dos factores (2FA)}

\begin{reglanegocio}[Reglas de 2FA]
\begin{itemize}
  \item El usuario activa 2FA generando un \texttt{two\_factor\_secret} (TOTP).
  \item Debe confirmar el código una vez; la fecha se registra en \texttt{two\_factor\_confirmed\_at}.
  \item Un usuario con 2FA no confirmado \textbf{no puede} acceder a canales de tiempo real (Broadcasting).
  \item La verificación de email (\texttt{email\_verified\_at}) es independiente del 2FA.
  \item El hash de verificación de email usa \texttt{sha256(email)}.
\end{itemize}
\end{reglanegocio}

\subsection{Identidades externas (OAuth)}

\begin{itemize}
  \item Los usuarios pueden vincular una identidad de Google via \texttt{auth.user\_identities}.
  \item Cada identidad almacena: \texttt{provider}, \texttt{provider\_uid}, \texttt{provider\_email}.
  \item Una identidad vinculada a otro usuario bloquea el inicio de sesión con ese proveedor (\texttt{isGoogleLinkedToAnotherUser}).
\end{itemize}


% ═══════════════════════════════════════════════════
\section{Incidencias — Ciclo de Vida}
% ═══════════════════════════════════════════════════

\subsection{Creación de incidencia}

\begin{reglanegocio}[Reglas de creación]
\begin{itemize}
  \item Requiere permiso \texttt{incidents.create} (Admin y Ciudadano).
  \item Campos obligatorios: \texttt{title}, \texttt{description}, \texttt{category\_id}, \texttt{territorial\_unit\_id}.
  \item La \texttt{priority\_id} \textbf{solo} puede ser definida por usuarios con rol de gestión (Admin, Supervisor). Si un ciudadano intenta enviarla, se ignora.
  \item Se genera un código único secuencial (\texttt{INC-YYYY-NNNNN}).
  \item El estado inicial es siempre \textbf{NUEVA}.
  \item Las coordenadas (\texttt{latitude}, \texttt{longitude}) se sincronizan automáticamente al campo PostGIS \texttt{location} via trigger de PostgreSQL.
  \item La fecha límite (\texttt{due\_date}) se calcula automáticamente a partir de las horas SLA de la prioridad via trigger.
\end{itemize}
\end{reglanegocio}

\subsection{Clasificación}

\begin{table}[H]
\centering
\begin{tabularx}{\textwidth}{l X l}
  \toprule
  \textbf{Categoría} & \textbf{Subcategorías} & \textbf{Color} \\
  \midrule
  Vialidad             & Bache, Semáforo dañado, Señalización vial, Hundimiento, Pavimento deteriorado & \textcolor[HTML]{6366F1}{$\blacksquare$} \\
  Servicios Públicos   & Fuga de agua, Alcantarilla tapada, Falta de agua, Drenaje colapsado & \textcolor[HTML]{0EA5E9}{$\blacksquare$} \\
  Alumbrado Público    & Luminaria apagada, Poste dañado, Cable caído, Zona sin iluminación & \textcolor[HTML]{F59E0B}{$\blacksquare$} \\
  Espacios Públicos    & Parque descuidado, Mobiliario dañado, Juegos infantiles rotos & \textcolor[HTML]{22C55E}{$\blacksquare$} \\
  Recolección de Residuos & Basura acumulada, Contenedor lleno, Residuos peligrosos & \textcolor[HTML]{A855F7}{$\blacksquare$} \\
  Seguridad            & Vandalismo, Zona insegura, Obstrucción de vía & \textcolor[HTML]{EF4444}{$\blacksquare$} \\
  \bottomrule
\end{tabularx}
\caption{Categorías y subcategorías de incidencias}
\end{table}

\subsection{Prioridades y SLA}

\begin{table}[H]
\centering
\begin{tabular}{l c c c c}
  \toprule
  \textbf{Prioridad} & \textbf{Nivel} & \textbf{SLA (horas)} & \textbf{Peso} & \textbf{Color} \\
  \midrule
  Crítica & 1 & 8   & 5 & \textcolor[HTML]{DC2626}{$\blacksquare$} \\
  Alta    & 2 & 24  & 3 & \textcolor[HTML]{F97316}{$\blacksquare$} \\
  Media   & 3 & 72  & 2 & \textcolor[HTML]{EAB308}{$\blacksquare$} \\
  Baja    & 4 & 168 & 1 & \textcolor[HTML]{22C55E}{$\blacksquare$} \\
  \bottomrule
\end{tabular}
\caption{Niveles de prioridad con SLA}
\end{table}

\begin{nota}[Peso de prioridad]
El \textbf{peso} se utiliza para calcular la carga de trabajo (\textit{workload points}) de un operador. Cada incidencia activa asignada suma su peso a la carga total del operador.
\end{nota}

\subsection{Máquina de estados}

Las incidencias transitan por un flujo de estados controlado por transiciones explícitas.

\begin{figure}[H]
\centering
\begin{tikzpicture}[
    node distance=2.4cm and 3.5cm,
    state/.style={
      rectangle, rounded corners=6pt, draw=#1!70, fill=#1!10,
      minimum width=2.6cm, minimum height=1cm,
      font=\small\bfseries, text=#1!90, align=center
    },
    trans/.style={-{Stealth[length=6pt]}, thick, color=primary!60},
    label/.style={font=\tiny, color=primary!50, fill=white, inner sep=1.5pt}
  ]

  \node[state=muted]   (nueva)      {NUEVA};
  \node[state=accent, right=of nueva]   (revision)   {EN\_REVISIÓN};
  \node[state=warning, right=of revision]  (progreso)   {EN\_PROGRESO};
  \node[state=success!80!black, below right=1.8cm and 1.7cm of progreso] (resuelta)   {RESUELTA};
  \node[state=success, below=2cm of resuelta]  (cerrada)    {CERRADA};
  \node[state=danger, below=2cm of nueva]   (rechazada)  {RECHAZADA};
  \node[state=warning!80!red, below=2cm of progreso]  (reabierta)  {REABIERTA};

  % Transiciones principales
  \draw[trans] (nueva)     -- node[label, above] {A,S} (revision);
  \draw[trans] (revision)  -- node[label, above] {A,S} (progreso);
  \draw[trans] (progreso)  -- node[label, above right] {A,S,O} (resuelta);
  \draw[trans] (resuelta)  -- node[label, right] {A,S} (cerrada);

  % Rechazos
  \draw[trans, danger!60] (nueva.south)     -- node[label, left] {A,S} (rechazada.north);
  \draw[trans, danger!60] (revision.south) to[bend right=20] node[label, left] {A,S} (rechazada.north east);

  % Reaperturas
  \draw[trans, warning!70!red] (cerrada.west)     -- node[label, below] {A,S} (reabierta.east);
  \draw[trans, warning!70!red] (rechazada.east)   -- node[label, below] {A,S} (reabierta.west);

  % Reabierta vuelve a revisión
  \draw[trans, accent!60] (reabierta.north) -- node[label, right] {A,S} (progreso.south);

\end{tikzpicture}
\caption{Diagrama de estados de una incidencia. A = Admin, S = Supervisor, O = Operador.}
\end{figure}

\begin{table}[H]
\centering
\small
\begin{tabularx}{\textwidth}{l l l c X}
  \toprule
  \textbf{Origen} & \textbf{Destino} & \textbf{Roles} & \textbf{Comentario} & \textbf{Descripción} \\
  \midrule
  NUEVA       & EN\_REVISIÓN  & A, S    & No  & Supervisor toma la incidencia para evaluación \\
  NUEVA       & RECHAZADA     & A, S    & Sí  & Se descarta la incidencia \\
  EN\_REVISIÓN & EN\_PROGRESO & A, S    & No  & Se asigna operador y comienza la resolución \\
  EN\_REVISIÓN & RECHAZADA    & A, S    & Sí  & Se descarta tras revisión \\
  EN\_PROGRESO & RESUELTA     & A, S, O & Sí  & El operador completa la resolución \\
  RESUELTA    & CERRADA       & A, S    & No  & El supervisor confirma resolución satisfactoria \\
  CERRADA     & REABIERTA     & A, S    & Sí  & Reapertura por resolución insatisfactoria \\
  RECHAZADA   & REABIERTA     & A, S    & Sí  & Se reconsidera el rechazo \\
  REABIERTA   & EN\_REVISIÓN  & A, S    & No  & Vuelve al flujo de revisión \\
  \bottomrule
\end{tabularx}
\caption{Transiciones válidas entre estados}
\end{table}

\begin{restriccion}[Reglas de transición]
\begin{itemize}
  \item Solo se permiten transiciones explícitamente definidas en \texttt{core.state\_transitions}.
  \item Cada transición verifica el rol del usuario autenticado contra \texttt{allowed\_roles}.
  \item Si \texttt{requires\_comment = true}, el comentario no puede ser nulo ni contener solo espacios.
  \item La reapertura requiere adicionalmente el permiso \texttt{incidents.reopen}.
  \item \textbf{RESUELTA no es un estado final}. Solo CERRADA y RECHAZADA son estados finales.
\end{itemize}
\end{restriccion}

\subsection{Campos temporales automáticos}

\begin{table}[H]
\centering
\begin{tabularx}{\textwidth}{l l X}
  \toprule
  \textbf{Campo} & \textbf{Se establece cuando} & \textbf{Lógica} \\
  \midrule
  \texttt{resolution\_date}           & Estado final (\texttt{is\_final\_state}) & Se asigna \texttt{now()} la primera vez; se preserva si ya existía \\
  \texttt{reopened\_at}               & Estado = REABIERTA  & Se registra el momento de reapertura \\
  \texttt{previous\_resolution\_date} & Estado = REABIERTA  & Se copia \texttt{resolution\_date} antes de anularlo \\
  \texttt{rejected\_at}               & Estado = RECHAZADA  & Se registra el momento del rechazo \\
  \texttt{due\_date}                  & Al asignar prioridad & Calculado por trigger PostgreSQL \\
  \bottomrule
\end{tabularx}
\caption{Campos temporales gestionados por el sistema}
\end{table}

\subsection{Solicitudes de cambio de estado (Operadores)}

\begin{reglanegocio}[Flujo de solicitud de cambio de estado]
\begin{enumerate}
  \item Un \textbf{operador} con asignación activa puede solicitar el cambio de estado de \texttt{EN\_PROGRESO} a \texttt{RESUELTA}.
  \item La solicitud se registra como \texttt{StateChangeRequest} con estado \texttt{pending}.
  \item Un \textbf{supervisor} o \textbf{administrador} aprueba o rechaza la solicitud.
  \item Si se aprueba, el sistema ejecuta el cambio de estado real.
  \item Solo puede existir una solicitud pendiente por incidencia a la vez.
  \item El operador debe tener una asignación activa en la incidencia.
\end{enumerate}
\end{reglanegocio}


% ═══════════════════════════════════════════════════
\section{Asignaciones y Estructura Operativa}
% ═══════════════════════════════════════════════════

\subsection{Asignación de operadores a incidencias}

\begin{reglanegocio}[Reglas de asignación]
\begin{itemize}
  \item Requiere permiso \texttt{incidents.assign} (Admin, Supervisor).
  \item La incidencia \textbf{debe tener una prioridad asignada} antes de asignar operadores.
  \item Soporta un operador \textbf{principal} (\texttt{primary}) y múltiples operadores de \textbf{apoyo} (\texttt{support}).
  \item Las asignaciones previas se desactivan y se crean las nuevas.
  \item El campo \texttt{current\_assigned\_id} de la incidencia se sincroniza con el operador principal.
  \item No se permite asignar operadores a incidencias en estado \texttt{CERRADA}.
\end{itemize}
\end{reglanegocio}

\subsection{Estructura de equipos}

\begin{itemize}
  \item \textbf{SupervisorProfile}: Perfil de supervisor vinculado a zonas territoriales.
  \item \textbf{OperatorProfile}: Perfil operativo con control de carga:
  \begin{itemize}
    \item \texttt{incident\_capacity}: Capacidad general del operador.
    \item \texttt{max\_active\_incidents}: Máximo de incidencias activas simultáneas (default: 10).
    \item \texttt{max\_workload\_points}: Puntos máximos de carga de trabajo (default: 20).
  \end{itemize}
  \item \textbf{SupervisorOperatorAssignment}: Relación supervisor--operador con fechas de asignación y estado activo.
  \item \textbf{UserTerritory}: Vinculación de usuarios a unidades territoriales con fechas y estado.
\end{itemize}


% ═══════════════════════════════════════════════════
\section{Jerarquía Territorial}
% ═══════════════════════════════════════════════════

\subsection{Modelo jerárquico}

El sistema implementa una jerarquía territorial basada en la estructura administrativa de Ecuador:

\begin{figure}[H]
\centering
\begin{tikzpicture}[
    level/.style={font=\small\bfseries, text=primary},
    every node/.style={draw=accent!40, fill=accent!5, rounded corners=3pt, minimum width=3cm, minimum height=0.8cm, font=\small},
    level distance=1.2cm, sibling distance=4cm,
    edge from parent/.style={draw=accent!50, -{Stealth[length=4pt]}, thick}
  ]
  \node {País (Ecuador)}
    child { node {Zona Operativa}
      child { node {Provincia}
        child { node {Cantón}
          child { node {Parroquia}
            child { node {Sector} }
          }
        }
      }
    };
\end{tikzpicture}
\caption{Jerarquía territorial}
\end{figure}

\begin{restriccion}[Reglas de validación jerárquica]
\begin{itemize}
  \item Cada tipo de unidad tiene un padre válido específico (ej: Provincia solo puede ser hija de Zona Operativa).
  \item No se permiten ciclos (un nodo no puede ser ancestro de sí mismo).
  \item La validación se realiza en el dominio via \texttt{TerritorialHierarchyRules}.
  \item Los datos base provienen del dataset CGE 2026 con 24 provincias, 221 cantones y 1024 parroquias.
  \item Las zonas operativas agrupan provincias según distribución geográfica (8 zonas).
\end{itemize}
\end{restriccion}

\subsection{Resolución de zona por coordenadas}

Cuando una incidencia tiene coordenadas geográficas, el sistema puede resolver la zona operativa correspondiente. Esto permite notificar al supervisor de la zona correcta aunque la incidencia haya sido registrada con una unidad territorial diferente.


% ═══════════════════════════════════════════════════
\section{Notificaciones}
% ═══════════════════════════════════════════════════

Las notificaciones internas del sistema se almacenan en \texttt{core.notifications} y se vinculan a la incidencia origen.

\subsection{Eventos que generan notificaciones}

\begin{table}[H]
\centering
\begin{tabularx}{\textwidth}{l l X}
  \toprule
  \textbf{Tipo} & \textbf{Destinatario} & \textbf{Evento} \\
  \midrule
  \texttt{STATUS\_CHANGE}     & Reportero, operadores asignados & Cambio de estado de una incidencia \\
  \texttt{ASSIGNMENT}         & Operador asignado               & Nueva asignación a una incidencia \\
  \texttt{NEW\_COMMENT}       & Reportero / Operadores          & Nuevo comentario o evidencia adjunta \\
  \texttt{STATUS\_CHANGE}     & Supervisor                      & Rechazo o cierre de incidencia crítica \\
  \texttt{STATE\_REQUEST}     & Supervisor                      & Solicitud de cambio de estado por operador \\
  \bottomrule
\end{tabularx}
\caption{Tipos de notificaciones}
\end{table}


% ═══════════════════════════════════════════════════
\section{Comentarios y Evidencia}
% ═══════════════════════════════════════════════════

\subsection{Comentarios}

\begin{itemize}
  \item Dos tipos: \textbf{públicos} (visibles para todos) e \textbf{internos} (solo equipo operativo).
  \item Para crear comentarios internos se requiere \texttt{comments.internal}.
  \item Un ciudadano \textbf{no puede} crear comentarios internos; si intenta marcar \texttt{is\_internal = true}, recibe un \texttt{403 Forbidden}.
  \item En las respuestas de detalle de incidencia, los comentarios internos se filtran automáticamente para usuarios sin el permiso correspondiente.
\end{itemize}

\subsection{Evidencia fotográfica}

\begin{itemize}
  \item Los adjuntos se almacenan en un sistema de archivos S3-compatible (RustFS).
  \item Formatos permitidos: \texttt{jpg}, \texttt{jpeg}, \texttt{png}.
  \item Tamaño máximo: 10 MB por archivo.
  \item Si el procesamiento falla, se notifica automáticamente al equipo administrativo.
\end{itemize}


% ═══════════════════════════════════════════════════
\section{Auditoría}
% ═══════════════════════════════════════════════════

El sistema implementa auditoría automática sobre todas las entidades que usan el trait \texttt{Auditable}.

\begin{reglanegocio}[Eventos auditados]
\begin{itemize}
  \item \textbf{created}: Snapshot completo de la nueva entidad.
  \item \textbf{updated}: Solo los campos que cambiaron (old/new).
  \item \textbf{deleted}: Snapshot de la entidad eliminada.
  \item \textbf{restored}: Snapshot tras restauración de soft delete.
\end{itemize}
Cada registro incluye: usuario autenticado, dirección IP, URL, user agent, y etiquetas de tabla.

Campos sensibles (\texttt{password}, \texttt{remember\_token}, \texttt{two\_factor\_secret}) se excluyen automáticamente de los snapshots.
\end{reglanegocio}

\begin{nota}[Entidades auditables]
Las entidades \texttt{User} e \texttt{Incident} implementan el trait \texttt{Auditable}, lo que asegura trazabilidad completa de todo cambio realizado sobre usuarios e incidencias del sistema.
\end{nota}


% ═══════════════════════════════════════════════════
\section{Reglas de Negocio — Resumen}
% ═══════════════════════════════════════════════════

\begin{enumerate}[label=\textbf{RN-\arabic*}, leftmargin=2.5cm]
  \item Un usuario tiene exactamente un rol.
  \item Solo Admin y Ciudadano pueden crear incidencias.
  \item Solo Admin y Supervisor pueden asignar prioridad a una incidencia.
  \item La prioridad es obligatoria antes de asignar operadores.
  \item Las transiciones de estado se validan contra la tabla de transiciones y los roles del usuario.
  \item Las transiciones que requieren comentario rechazan comentarios vacíos o de solo espacios.
  \item RESUELTA no es estado final; solo CERRADA y RECHAZADA son finales.
  \item La reapertura requiere el permiso \texttt{incidents.reopen} y un comentario.
  \item La reapertura limpia la fecha de resolución y preserva la anterior.
  \item Los operadores solo pueden solicitar cambios de estado, no ejecutarlos directamente (excepto EN\_PROGRESO → RESUELTA).
  \item Los comentarios internos solo son visibles para roles con \texttt{comments.internal}.
  \item Un usuario con 2FA no confirmado no puede autorizar canales de Broadcasting.
  \item Las coordenadas se sincronizan al campo PostGIS via trigger.
  \item La fecha límite se calcula automáticamente a partir de las horas SLA.
  \item El cierre de una incidencia desactiva todas las asignaciones activas.
  \item La carga de trabajo del operador se controla por puntos (peso de prioridad) y máximo de incidencias activas.
  \item Todo cambio sobre usuarios e incidencias genera un registro de auditoría automático.
  \item Las identidades OAuth vinculadas a otro usuario bloquean el inicio de sesión con ese proveedor.
\end{enumerate}


% ═══════════════════════════════════════════════════
\section{Infraestructura Técnica}
% ═══════════════════════════════════════════════════

\begin{table}[H]
\centering
\begin{tabularx}{\textwidth}{l X}
  \toprule
  \textbf{Componente} & \textbf{Tecnología} \\
  \midrule
  Backend       & Laravel 12, PHP 8.3 \\
  Base de datos & PostgreSQL 16 con PostGIS \\
  Caché/Sesión  & Redis \\
  Archivos      & RustFS (S3-compatible) \\
  Email (dev)   & Mailpit \\
  Frontend      & HTML5 + CSS3 + JavaScript vanilla + Bootstrap 5 + AdminLTE 3 \\
  API           & REST stateless con tokens Sanctum \\
  Tests         & PHPUnit (Unit + Feature) \\
  \bottomrule
\end{tabularx}
\caption{Stack tecnológico}
\end{table}

\vfill
\begin{center}
  \textcolor{muted}{\rule{0.4\textwidth}{0.5pt}}\\[0.5em]
  \small\textcolor{muted}{SGI — Sistema de Gestión de Incidencias Georreferenciadas · Julio 2026}
\end{center}

\end{document}
